% ══════════════════════════════════════════════════════════════════════════════
%  Chapter 19: Wallet UI (Tauri)
% ══════════════════════════════════════════════════════════════════════════════
\chapter{Wallet UI (Tauri)}
\label{ch:wallet-tauri}

\begin{learnbox}
\begin{itemize}
  \item Build a user-facing wallet interface using Tauri 2 (Rust backend \texttt{+} React frontend)
  \item Implement secure persistent storage for wallet keys and transaction history
  \item Design a modern React UI for real wallet workflows with transaction signing
  \item Integrate Tauri commands for database operations and transaction broadcast
\end{itemize}
\end{learnbox}

\lettrine{T}{his} chapter explains the \code{\index{bitcoin-wallet-ui-tauri}bitcoin-wallet-ui-tauri} \index{Cargo!crate}crate---the same \index{wallet}wallet \index{functionality}functionality as Chapter~\ref{ch:wallet-iced}, but built with \index{Tauri}Tauri 2 instead of pure \index{Iced}Iced.

\begin{notebox}[Architecture]
Like the \index{admin dashboard}admin UI \index{comparison}comparison (Chapters~\ref{ch:desktop-iced} and~\ref{ch:desktop-tauri}), this \index{wallet}wallet demonstrates the same \textbf{\index{feature}feature set} with two different \index{architecture}architectures. Chapter~\ref{ch:wallet-iced} is pure \index{Rust!backend}Rust (\index{Iced}Iced \index{MVU}MVU); this chapter is \index{Tauri}Tauri (\index{Rust!backend}Rust \index{backend}backend \texttt{+} \index{React}React \index{frontend}frontend). If you read Chapter~\ref{ch:wallet-iced} first, you will see where the \index{UI}UI \index{logic}logic diverges and how the same \index{blockchain}blockchain \index{operation}operations are exposed through different \index{framework}frameworks.
\end{notebox}

The \index{Tauri}Tauri \index{wallet}wallet adds \textbf{\index{persistent}persistent \index{encryption}encrypted \index{storage}storage} of \index{private key}private \index{key}keys and \index{transaction}transaction \index{history}history. Unlike the \index{admin dashboard}admin UI (which connects to a remote \index{blockchain}blockchain \index{node}node), the \index{wallet}wallet UI manages \index{key}keys locally and signs \index{transaction}transactions in the \index{Rust!backend}Rust \index{backend}backend before \index{broadcast}broadcasting them.

\section{What This UI Is (In One Sentence)}
\label{sec:ch19-oneline}

The Tauri wallet UI is a \textbf{two-language user-facing application} that adds \textbf{secure key storage}: a Rust backend exposes 18 \code{\#[tauri::command]} handlers for wallet operations (address generation, balance queries, transaction signing and broadcast), while a React frontend provides a modern dashboard for managing addresses, sending transactions, and viewing history.

\section{Architectural Differences: Iced vs. Tauri Wallet}
\label{sec:ch19-comparison}

Both the Iced (Chapter~\ref{ch:wallet-iced}) and Tauri (this chapter) wallet implementations expose the same operations:

\begin{table}[htbp]
\caption{Wallet UI: Iced vs. Tauri}
\label{tab:wallet-iced-vs-tauri}
\centering\small
\begin{tabularx}{0.95\textwidth}{lXX}
\toprule
\textbf{Aspect} & \textbf{Iced Wallet (Ch. 18)} & \textbf{Tauri Wallet (Ch. 19)} \\
\midrule
Languages & Rust only & Rust backend \texttt{+} React/TypeScript \\
State persistence & SQLite with encrypted password & SQLite via Rust backend, React state (Zustand) \\
Key storage & In-memory + database & In Rust backend (\code{State<RwLock>}), never exposed to JS \\
Transaction signing & Rust (sync) & Rust backend via command, never touches frontend \\
Event handling & \code{Message} enum & React Router \texttt{+} React Query \\
UI styling & Iced widgets (programmatic) & Tailwind CSS \texttt{+} React components \\
\bottomrule
\end{tabularx}
\end{table}

\begin{infobox}[Security Note]
Both implementations follow the same security principle: \textbf{private keys stay in Rust}. In Iced, keys are in memory and encrypted at rest in SQLite. In Tauri, keys are in Rust backend state and never serialized to JSON/JavaScript. Transaction signing happens in Rust; the frontend only receives the signed transaction hash.
\end{infobox}

% ──────────────────────────────────────────────────────────────────────────────
\section{Rust Backend Architecture}
\label{sec:ch19-rust-backend}

\begin{listing}[H]
\caption{Tauri wallet main (\code{src-tauri/src/main.rs})}
\label{lst:ch19-main}
\begin{minted}[fontsize=\footnotesize]{rust}
use tauri::State;
use std::sync::RwLock;

#[derive(Clone, serde::Serialize, serde::Deserialize)]
pub struct WalletState {
    pub current_wallet: Option<String>,
    pub wallets: Vec<String>,
    pub database_password: String,
}

#[derive(Clone)]
pub struct AppState {
    pub wallet: RwLock<WalletState>,
    pub base_url: String,
}

fn main() {
    let app_state = AppState {
        wallet: RwLock::new(WalletState {
            current_wallet: None,
            wallets: Vec::new(),
            database_password: generate_database_password(),
        }),
        base_url: "http://127.0.0.1:8080".to_string(),
    };

    tauri::Builder::default()
        .manage(app_state)
        .invoke_handler(tauri::generate_handler![
            cmd_list_wallets,
            cmd_create_wallet,
            cmd_select_wallet,
            cmd_get_balance,
            cmd_get_addresses,
            cmd_generate_address,
            cmd_get_transaction_history,
            cmd_get_unspent_outputs,
            cmd_send_transaction,
            cmd_sign_transaction,
            cmd_get_tx_status,
            cmd_export_wallet,
            cmd_import_wallet,
            cmd_backup_wallet,
            cmd_set_passphrase,
            cmd_unlock_wallet,
            cmd_get_wallet_info,
            cmd_cancel_transaction,
        ])
        .run(tauri::generate_context!())
        .expect("error while running tauri application");
}

fn generate_database_password() -> String {
    use std::collections::hash_map::DefaultHasher;
    use std::hash::{Hash, Hasher};

    let mut hasher = DefaultHasher::new();
    if let Ok(username) = std::env::var("USER") {
        username.hash(&mut hasher);
    }
    if let Some(home) = dirs::home_dir() {
        home.to_string_lossy().hash(&mut hasher);
    }
    format!("{:x}", hasher.finish())
}
\end{minted}
\end{listing}

\textbf{Walkthrough: the command handler pattern:}

Every command follows the same structure:

\begin{enumerate}
  \item Extract \code{State<AppState>} parameter injected by Tauri.
  \item Call a corresponding \code{WalletClient} method from the bitcoin-api crate.
  \item Return the result (or error) as JSON over IPC.
\end{enumerate}

Crucially, transaction signing and key management happen in the Rust backend---the frontend never touches private keys.

% ──────────────────────────────────────────────────────────────────────────────
\section{React Frontend: State Management and Commands}
\label{sec:ch19-frontend}

\begin{listing}[H]
\caption{React Query hooks (\code{src/hooks/useWallet.ts})}
\label{lst:ch19-hooks}
\begin{minted}[fontsize=\footnotesize]{typescript}
import { useQuery, useMutation } from '@tanstack/react-query';
import * as commands from '../commands';

export function useWalletList() {
    return useQuery({
        queryKey: ['wallets'],
        queryFn: () => commands.listWallets(),
        staleTime: 30000,
    });
}

export function useBalance(address: string | null) {
    return useQuery({
        queryKey: ['balance', address],
        queryFn: () => address ? commands.getBalance(address) : Promise.resolve(0),
        enabled: !!address,
        staleTime: 10000,
    });
}

export function useSendTransaction() {
    return useMutation({
        mutationFn: (params: { from: string; to: string; amount: number }) =>
            commands.sendTransaction(params.from, params.to, params.amount),
    });
}

export function useCreateWallet() {
    return useMutation({
        mutationFn: (label?: string) => commands.createWallet(label),
    });
}
\end{minted}
\end{listing}

React Query manages server state (data from the Bitcoin API), while local component state handles UI (forms, modals, etc.):

\begin{enumerate}
  \item \textbf{useQuery}: automatic caching and refetching for reads (balance, transaction history).
  \item \textbf{useMutation}: explicit side effects for writes (create wallet, send transaction).
\end{enumerate}

Each query/mutation calls the corresponding function in \code{commands.ts}, which uses Tauri's \code{invoke()}.

% ──────────────────────────────────────────────────────────────────────────────
\section{Persistent Database Layer}
\label{sec:ch19-database}

Both Iced (Chapter~\ref{ch:wallet-iced}) and Tauri (this chapter) wallets use the same encrypted SQLite database (Chapter~\ref{ch:embedded-db}) for persistence. The difference is where the encryption happens:

\begin{itemize}
  \item \textbf{Iced}: database initialization happens in \code{main.rs} before MVU starts.
  \item \textbf{Tauri}: database initialization happens in a startup command; the frontend can choose when to unlock it (e.g., with a passphrase).
\end{itemize}

\begin{listing}[htbp]
\caption{Database initialization command (\code{src-tauri/src/commands/database.rs})}
\label{lst:ch19-db-cmd}
\begin{minted}[fontsize=\footnotesize]{rust}
use crate::AppState;
use tauri::State;

#[tauri::command]
pub async fn cmd_init_database(
    state: State<'_, AppState>,
    password: Option<String>,
) -> Result<(), String> {
    let wallet_state = state.wallet.read()
        .map_err(|e| format!("Failed to read state: {}", e))?;

    // Use provided password or deterministic password
    let pwd = password.unwrap_or_else(|| {
        wallet_state.database_password.clone()
    });

    // Initialize SQLite with encryption (XChaCha20 via sqlcipher or similar)
    database::init_with_password(&pwd)
        .await
        .map_err(|e| e.to_string())
}

#[tauri::command]
pub async fn cmd_save_wallet(
    wallet_address: String,
    wallet_data: serde_json::Value,
) -> Result<(), String> {
    database::save_wallet(&wallet_address, &wallet_data)
        .await
        .map_err(|e| e.to_string())
}

#[tauri::command]
pub async fn cmd_load_wallet(
    wallet_address: String,
) -> Result<serde_json::Value, String> {
    database::load_wallet(&wallet_address)
        .await
        .map_err(|e| e.to_string())
}
\end{minted}
\end{listing}

The database persists:
\begin{itemize}
  \item Wallet metadata (addresses, labels).
  \item Transaction history (cached from the blockchain node).
  \item Encrypted key material (never directly accessed by the frontend).
\end{itemize}

% ──────────────────────────────────────────────────────────────────────────────
\begin{chaptersummary}
\item We built a user-facing Tauri wallet with a Rust backend and React frontend, mirroring the Iced wallet from Chapter~\ref{ch:wallet-iced}.
\item Command handlers expose 18 wallet operations (create, balance, send, history, etc.), all with transaction signing in Rust.
\item React Query manages server state (balances, transactions) and integrates with Tauri's IPC via \code{invoke()}.
\item The encrypted database persists wallet metadata and transaction history (Chapter~\ref{ch:embedded-db} covers the database itself).
\item Security is maintained by keeping private keys in Rust backend state---they never touch the JavaScript/React frontend.
\end{chaptersummary}

% ──────────────────────────────────────────────────────────────────────────────
\section{Exercises}
\label{sec:ch19-exercises}

\begin{enumerate}
  \item \textbf{Add a Transaction Confirmation Dialog} --- Modify the ``Send'' command to return a transaction preview (recipients, amounts, fees) to the frontend. Add a React modal that displays this preview and requires explicit confirmation before broadcasting. Trace the IPC round-trip: React \(\to\) Tauri command \(\to\) Bitcoin API \(\to\) JSON response \(\to\) modal display.

  \item \textbf{Implement Wallet Recovery} --- Add a ``Restore Wallet'' feature: allow the user to input a seed phrase, call a backend command to derive addresses from the seed, and persist them to the database. Compare this with the ``Create Wallet'' flow.
\end{enumerate}

% ──────────────────────────────────────────────────────────────────────────────
\section{Further Reading}
\label{sec:ch19-reading}

\begin{itemize}
  \item \textbf{Tauri Documentation} --- API reference and examples at \url{https://tauri.app/docs/}.
  \item \textbf{React Query Documentation} --- Server state management at \url{https://tanstack.com/query/}.
  \item \textbf{Chapter~\ref{ch:wallet}: Wallet System} --- The wallet implementation that both UIs consume.
  \item \textbf{Chapter~\ref{ch:embedded-db}: Embedded Database} --- SQLite encryption and persistence layer.
\end{itemize}

